\chapter{Electric charge and field}\label{ch:efield}

Part~II trades mechanics for fields, and the first one is the cleanest:
Coulomb's $1/r^{2}$ law looks like gravitation with signs --- but the
signs change everything (charges attract \emph{and} repel, so fields can
cancel).

\section{Charge: quantised, conserved, mobile in metal}\label{sec:ef-charge}

Charge comes in units of $e=1.602\times10^{-19}$\,C, is conserved in
every process, and lives in two kinds of material
(figure~\ref{fig:cond}): \emph{conductors} let electrons move freely
(charge spreads to the surface) while \emph{insulators} pin it where you
put it. Induction and polarization follow from mobility alone.

\begin{figure}[htbp]\centering
\begin{tikzpicture}[scale=0.9,font=\scriptsize]
\draw[very thick,rounded corners] (-2.5,-1) rectangle (-0.5,1);
\node at (-1.5,0.6) {conductor};
\foreach \x in {-2.2,-1.8,-1.4,-1.0,-0.6}{\node[amber] at (\x,-0.8) {$+$};}
\node at (-1.5,-1.5) {$E$ inside $=0$; charge lives on skin};
\draw[very thick,rounded corners] (0.8,-1) rectangle (2.8,1);
\node at (1.8,0.6) {insulator};
\foreach \x/\y in {1.1/-0.5, 1.5/0.2, 1.9/-0.3, 2.3/0.4, 2.5/-0.6}{\node[amber] at (\x,\y) {$+$};}
\node at (1.8,-1.5) {charge stays where rubbed};
\end{tikzpicture}

\caption{The material divide: conductors send charge to the surface and
kill the interior field; insulators keep it where deposited.}
\label{fig:cond}
\end{figure}

\section{Coulomb's law --- superposition is the real law}\label{sec:ef-coul}

\begin{equation}
  \mathbf F=k\frac{q_1q_2}{r^{2}}\hat r,\qquad
  k=\frac{1}{4\pi\varepsilon_0}=8.99\times10^{9}\,\mathrm{Nm^{2}/C^{2}},
  \label{eq:coul}
\end{equation}
and --- the loaded part --- \emph{forces add}. The force on a charge is
the vector sum over every other charge, one pair at a time, with no
three-body terms. That linearity is what makes field theory possible.

\begin{example}[Superposition, vectors not scalars]\label{ex:superp}
Charges $+2$\,\textmu C at the origin and $-1$\,\textmu C at
$(3,0)$\,cm. Find the force on $+1$\,\textmu C at $(0,4)$\,cm.
\par\smallskip
\textbf{Strategy.} Compute each force as a vector and add.
\begin{solution}
$r_1=0.04$\,m along $\hat y$, so
\[\mathbf F_1=\frac{k(2\mu)(1\mu)}{0.0016}\,\hat y=11.24\,\hat y\ \text{N}.\]
$r_2=\sqrt{0.03^{2}+0.04^{2}}=0.05$\,m; the negative charge pulls along
$(3,-4)/5$, giving
\[\mathbf F_2=\frac{k(1\mu)(1\mu)}{0.0025}\cdot\frac{(3,-4)}{5}
=(2.16,-2.88)\ \text{N}.\]
Sum: $(2.16,\,8.36)$\,N, magnitude $8.6$\,N at $75^{\circ}$ ---
not the $14.8$\,N of scalar addition. \checkmark
\end{solution}
\end{example}

\section{The field: force per test charge}\label{sec:ef-field}

$\mathbf E=\mathbf F/q$ --- defined so the source needn't care what test
charge arrives. A point charge's is
$\mathbf E=kq\hat r/r^{2}$ (figure~\ref{fig:pt}); a distribution's is the
vector sum over all its $dq$.

\begin{figure}[htbp]\centering
\begin{tikzpicture}[scale=0.95]
\foreach \ang in {0,45,...,315}{\draw[->,thick,accent] (0,0) -- ({\ang}:1.9);}
\foreach \ang in {0,45,...,315}{\draw[->,accent] (0,0) -- ({\ang}:2.6);}
\fill[amber] (0,0) circle (0.07) node[below,font=\scriptsize]{$+q$};
\node[font=\scriptsize] at (0,-3.0) {point charge: radial, density $\propto 1/r^{2}$};
\end{tikzpicture}

\caption{The point-charge field: radial, with line density standing in
for the $1/r^{2}$ fall-off.}
\label{fig:pt}
\end{figure}

\section{Field lines: reading, rules}\label{sec:ef-lines}

Lines start on $+$, end on $-$, never cross (the field is single-valued),
and their density reads magnitude (figure~\ref{fig:lines}). They are a
map, not a mechanism --- the physics is \eqref{eq:coul}.

\begin{figure}[htbp]\centering
\begin{tikzpicture}[scale=0.85,font=\scriptsize]
\fill[amber] (-1.6,0) circle (0.08) node[below]{$+q$};
\fill[amber] (1.6,0) circle (0.08) node[below]{$+q$};
\draw[->,accent] (-1.6,0) -- (-0.4,0.9);
\draw[->,accent] (1.6,0) -- (0.4,0.9);
\draw[->,accent] (-1.6,0) -- (-0.4,-0.9);
\draw[->,accent] (1.6,0) -- (0.4,-0.9);
\draw[->,accent] (-1.6,0) -- (-2.9,0);
\draw[->,accent] (1.6,0) -- (2.9,0);
\draw[->,accent] (-1.6,0) -- (-1.0,1.3);
\draw[->,accent] (1.6,0) -- (1.0,1.3);
\draw[->,accent] (-1.6,0) -- (-1.0,-1.3);
\draw[->,accent] (1.6,0) -- (1.0,-1.3);
\node at (0,-1.9) {two like charges: lines never meet, never cross};
\end{tikzpicture}

\caption{Two like charges: lines repel into a saddle. Between two equal
charges the field vanishes exactly at the midpoint.}
\label{fig:lines}
\end{figure}

\begin{figure}[htbp]\centering
\begin{tikzpicture}[scale=0.85]
\foreach \sgn/\xx in {+/-1.3, -/1.3}{
  \foreach \a in {60,120,240,300}{
    \draw[->,accent] (\xx,0) -- ++(\a:1.1);}}
\fill[amber] (-1.3,0) circle (0.08) node[below,font=\scriptsize]{$+q$};
\fill[accent] (1.3,0) circle (0.08) node[below,font=\scriptsize]{$-q$};
\draw[thick,->] (-0.9,0) -- (0.9,0);
\draw[thick,->] (-0.85,0.35) .. controls (0,0.9) .. (0.85,0.35);
\draw[thick,->] (-0.85,-0.35) .. controls (0,-0.9) .. (0.85,-0.35);
\draw[thick] (-1.0,0.6) .. controls (-0.3,1.4) and (0.3,1.4) .. (1.0,0.6);
\draw[thick] (-1.0,-0.6) .. controls (-0.3,-1.4) and (0.3,-1.4) .. (1.0,-0.6);
\node[font=\scriptsize] at (0,-1.8) {dipole: lines $+\to-$, far field $\propto p/r^{3}$};
\end{tikzpicture}

\caption{The dipole: far field falls as $p/r^{3}$ --- the first
\emph{non}-inverse-square pattern in the book.}
\label{fig:dip}
\end{figure}

\begin{example}[Dipole far field]\label{ex:dipole}
Show that on the dipole axis, $|E|\to2kp/r^{3}$ for $r\gg a$,
$p=qa$.
\par\smallskip
\textbf{Strategy.} Superpose two point-charge fields at distances
$r-a/2$ and $r+a/2$, then expand.
\begin{solution}
$E=kq[(r-a/2)^{-2}-(r+a/2)^{-2}]
=kq\frac{(r+a/2)^{2}-(r-a/2)^{2}}{(r^{2}-a^{2}/4)^{2}}
=kq\frac{2ra}{(r^{2})^{2}}\to\frac{2kp}{r^{3}}$.
The $1/r^{2}$ terms cancel; the \emph{difference} survives ---
the code check below finds the same $r^{-3}$. \checkmark
\end{solution}
\end{example}

\section{Continuous distributions: pay with integrals}\label{sec:ef-cont}

For charge spread over a line, ring or disk, integrate
$d\mathbf E=k\,dq\,\hat r/r^{2}$ element by element --- and let symmetry
discard every component that cancels
(figure~\ref{fig:ring}).

\begin{figure}[htbp]\centering
\begin{tikzpicture}[scale=1.1]
\draw[very thick,accent] (0,0) circle (1.1);
\foreach \i in {0,1,...,7}{\fill (1.1*\i*45:0) ;}
\foreach \ang in {0,45,...,315}{\fill ({1.1*cos(\ang)},{1.1*sin(\ang)}) circle (0.03);}
\coordinate (P) at (3.0,0);
\draw[dashed] (0,0) -- (P) node[right,font=\scriptsize]{$P$: on axis};
\draw[->,thick,amber] (0.95,0.57) -- (P);
\draw[->,thick,amber] (-0.95,-0.57) -- (P);
\node[font=\scriptsize,align=left] at (0,-1.6) {each $dq$ contributes $dE$;\\
transverse parts cancel by pairs};

\end{tikzpicture}

\caption{The ring's axis field: each $dq$'s transverse piece cancels its
opposite's --- only $dE_x$ survives.}
\label{fig:ring}
\end{figure}

\begin{example}[Field on the axis of a charged ring]\label{ex:ring}
A ring of radius $R$ carries $Q$ uniformly. Find $\mathbf E$ on the
axis at distance $x$.
\par\smallskip
\textbf{Strategy.} Every element is the same distance
$\sqrt{x^{2}+R^{2}}$; transverse components cancel.
\begin{solution}
Each $dq$ contributes axial
$dE_x=k\,dq\,x/(x^{2}+R^{2})^{3/2}$; integrating over the ring just
multiplies by $Q$:
$E=\dfrac{kQx}{(x^{2}+R^{2})^{3/2}}$.
Limits: $x\to0$ gives $0$ (symmetry wins at the centre), $x\gg R$ gives
$kQ/x^{2}$ (the ring becomes a point). The code confirms it to $0.00\%$.
\checkmark
\end{solution}
\end{example}

\begin{exercise}[easy]
Where on the axis is the ring's field maximum? (Differentiate.) (AP-C)
\end{exercise}

\begin{example}[Charged rod at its end]\label{ex:rod}
A rod of length $L$ carries $Q$ uniformly. Find $\mathbf E$ a distance
$d$ past its near end.
\par\smallskip
\textbf{Strategy.} Integrate element by element --- no symmetry to
help this time.
\begin{solution}
With $dq=\lambda\,dx$, $\lambda=Q/L$, at distance $r=d+L-x$ from the
element:
\[E=\int_0^{L}\frac{k\lambda\,dx}{(d+L-x)^{2}}
=k\lambda\left[\frac{1}{d}-\frac{1}{d+L}\right]
=\frac{kQ}{d(d+L)}.\]
Limits: $d\gg L$ recovers $kQ/d^{2}$; $d\to0$ diverges --- the rod
becomes its own endpoint's infinity. \checkmark
\end{solution}
\end{example}

\section{Motion in a uniform field: kinematics returns}\label{sec:ef-motion}

Inside a parallel-plate capacitor $E$ is constant, so a charge's
trajectory is chapter~1's projectile motion with $a=qE/m$
(figure~\ref{fig:defl}) --- oscilloscopes and inkjets are this
application sold as hardware.

\begin{figure}[htbp]\centering
\begin{tikzpicture}[scale=1.0]
\draw[thick] (0,0.6) -- (4.4,0.6);
\draw[thick] (0,-0.6) -- (4.4,-0.6);
\node[font=\scriptsize] at (4.7,0.6) {$+$};
\node[font=\scriptsize] at (4.7,-0.6) {$-$};
\draw[->,thick,amber] (-0.4,0) -- (0.5,0) node[below,font=\scriptsize]{$v_0$};
\draw[very thick,accent,domain=0:4.4,samples=40] plot (\x,{-0.05*\x*\x});
\draw[dashed] (4.4,-0.97) -- (4.4,-1.3) node[below,font=\scriptsize]{deflection $\propto E L^{2}/v_0^{2}$};
\node[font=\scriptsize] at (1.2,1.2) {constant $a_y$: projectile again};
\end{tikzpicture}

\caption{Electron deflection between plates: the projectile equation
wearing an electric costume.}
\label{fig:defl}
\end{figure}

\begin{example}[Oscilloscope deflection]\label{ex:deflect}
An electron ($v_0=2\times10^{7}$\,m/s) crosses plates of length
$L=4$\,cm with $E=2000$\,N/C. Find its deflection.
\par\smallskip
\textbf{Strategy.} Time in plates $t=L/v_0$; $y=\tfrac12at^{2}$.
\begin{solution}
$a=eE/m=1.6\times10^{-19}\times2000/9.11\times10^{-31}=3.5\times10^{14}$\,m/s$^{2}$;
$t=2\times10^{-9}$\,s;
$y=\tfrac12(3.5\times10^{14})(4\times10^{-18})=0.70$\,mm ---
and $\theta=\arctan(v_y/v_0)\approx10^{\circ}$ after exit. \checkmark
\end{solution}
\end{example}

\begin{lstlisting}[language=Python,caption={\texttt{code/ch09\_field\_lines.py} --- superposition check: dipole and ring vs theory.},label={lst:field}]
def E(charges, x, y, z=0.0):
    for q, xq, yq, zq in charges:
        dx, dy, dz = x-xq, y-yq, z-zq
        r3 = (dx*dx + dy*dy + dz*dz)**1.5
        ex += q*dx/r3; ey += q*dy/r3; ez += q*dz/r3
\end{lstlisting}

Output of a real run:
\begin{lstlisting}[numbers=none]
Dipole far-field ON axis (expect |E| -> 2p/r^3, p=qa=0.1):
  r=1.0: |E|=0.2010  theory=0.2000  ratio=1.005
  r=2.0: |E|=0.0250  theory=0.0250  ratio=1.001
  r=4.0: |E|=0.0031  theory=0.0031  ratio=1.000
Dipole EQUATOR (expect |E| -> p/r^3):
  r=1.0: |Ex|=0.0996  theory=0.1000
  r=2.0: |Ex|=0.0125  theory=0.0125

Ring N=200 on axis (expect Ex=x/(x^2+R^2)^1.5, Ey=Ez=0):
  x=0.2: Ex=1.2807  theory=1.2807  err=0.00%  trans=4.3e-16
  x=0.5: Ex=1.4142  theory=1.4142  err=0.00%  trans=1.0e-16
  x=1.0: Ex=0.7155  theory=0.7155  err=0.00%  trans=4.9e-18
  x=2.0: Ex=0.2283  theory=0.2283  err=0.00%  trans=5.2e-18
\end{lstlisting}

Naked superposition with $N=200$ ring elements nails every limit: the
dipole ratio converges to $1.000$ as $r/a$ grows, the ring matches the
analytic $x(x^{2}+R^{2})^{-3/2}$ to $0.00\%$, and the transverse
components --- $10^{-16}$ --- are symmetry made numeric.

\section{In engineering practice: electrostatics as a design spec}\label{sec:ef-eng}

Electrostatic precipitators run \S\ref{sec:ef-motion} at industrial
scale (corona-charge the soot, plate-field it out); inkjet heads steer
charged droplets by the deflection you just computed; and every high-voltage
lab's ``corona ring'' is a real conductor shaped to flatten the field
of \S\ref{sec:ef-lines} --- rounded because sharp corners concentrate
lines into breakdown.

\section{Common mistakes}\label{sec:ef-mistakes}

\begin{tripwire}[Field vs force conflated]
$\mathbf E$ is a property of the source; $\mathbf F=q\mathbf E$ needs
the test charge. ``Force at a point'' without a charge there is
meaningless --- but the field isn't.
\end{tripwire}
\begin{tripwire}[Scalar superposition]
$|E|=|E_1|+|E_2|$ only when parallel. Two perpendicular equal fields
give $\sqrt2\,|E_1|$, not $2|E_1|$.
\end{tripwire}
\begin{tripwire}[$dq$ element errors]
For a line, $dq=\lambda\,dx$ --- but on an arc, $dq=\lambda R\,d\theta$
and the distance to the field point is \emph{not} constant unless you
integrate.
\end{tripwire}
\begin{tripwire}[Sign of the test charge]
Negative charges feel $\mathbf F$ \emph{opposite} $\mathbf E$; tracing
an electron's path down field lines is backwards.
\end{tripwire}
\begin{tripwire}[Field ``used up'' by the probe]
A test charge measures without disturbing --- in the limit $q\to0$.
Real probes do disturb; that's a lab problem, not the definition's.
\end{tripwire}

\section{Summary}\label{sec:ef-summary}

\begin{table}[ht]\centering
\begin{tabularx}{\linewidth}{@{}lY@{}}
\toprule
\tablehead\thead{Idea} & \thead{One-line version}\\
\midrule
charge & quantised at $e$; mobile in conductor, pinned in insulator \\
Coulomb & $F=kq_1q_2/r^{2}$; forces add as vectors \\
field & $\mathbf E=\mathbf F/q$; point $kq\hat r/r^{2}$ \\
lines & start $+$/end $-$; density $\propto |E|$ \\
ring axis & $E=kQx/(x^{2}+R^{2})^{3/2}$ \\
dipole & far field $2kp/r^{3}$ (axis), $kp/r^{3}$ (equator) \\
uniform $E$ & projectile kinematics, $a=qE/m$ \\
\bottomrule
\end{tabularx}
\caption{Electric-field cheat sheet.}
\label{tab:efsummary}
\end{table}

\begin{wins}[Definition of done]
You can now: superpose Coulomb fields as vectors, read field-line maps,
integrate a continuous distribution by symmetry, derive ring/dipole
limits, and run charge-in-uniform-$E$ kinematics.
\end{wins}

\chapternotes{%
The field-first order follows MIT 8.02SC \cite{ocw801} modules 1--4 and
AP Physics C E\&M unit 8 \cite{apced}; Purcell--Morin's ch.~1 is the
authoritative charge/field treatment \cite{hrw}. The numeric checks are
the listing's real output.}

\begin{rulebox}[{Watch instead}]
\begin{enumerate}
\item \emph{Electric Charges and Forces --- Coulomb's Law --- Polarization}
--- Walter Lewin, MIT 8.02x Lecture 1, 47:13,
\url{https://youtu.be/x1-SibwIPM4} ---
\S\ref{sec:ef-charge}--\S\ref{sec:ef-coul} at full length (accessed
2026-10-11).
\item \emph{Electric field definition} --- Khan Academy, 13:46,
\url{https://youtu.be/_4VC3IHbuW8} --- \S\ref{sec:ef-field} worked
quickly (accessed 2026-10-11).
\end{enumerate}
\end{rulebox}

\section*{Further reading}
\addcontentsline{toc}{section}{Further reading}

\begin{itemize}
\item MIT OpenCourseWare, \emph{8.02SC Electricity and Magnetism},
modules 1--4 --- charge, field and continuous distributions. At
\url{https://ocw.mit.edu}, course ``8.02SC'' (verified 2026-10-11).
\item OpenStax, \emph{College Physics 2e}, ch.~18 ``Electric Charge and
Electric Field'' --- free text at \url{https://openstax.org} (verified
2026-10-11).
\item College Board \emph{AP Physics C (E\&M) CED}, Unit 8 ---
the exam blueprint for this chapter, at
\url{https://apstudents.collegeboard.org} (verified 2026-10-11).
\item E.~M. Purcell and D.~J. Morin, \emph{Electricity and Magnetism},
ch.~1 --- the deeper electrostatics; ISBN 978-1107014022.
\item R.~P. Feynman, \emph{Lectures on Physics}, vol.~II, ch.~1 ---
electromagnetism's stakes set in one chapter; free at Caltech's site.
\item Khan Academy, \emph{Electric charge, field, and potential} ---
drills matching \S\ref{sec:ef-coul}--\S\ref{sec:ef-cont}, at
\url{https://www.khanacademy.org} (verified 2026-10-11).
\end{itemize}

\section*{Additional exercises}
\addcontentsline{toc}{section}{Additional exercises}

\begin{exercise}[easy]
Two charges $+3$\,\textmu C and $-1$\,\textmu C sit $10$\,cm apart.
Find the force on a $+1$\,\textmu C test at the midpoint. (AP-C)
\end{exercise}
\begin{exercise}[easy]
A field points $+\hat x$ at a point: what does a proton do there ---
and an electron? (AP-C)
\end{exercise}
\begin{exercise}[easy]
Express $\mathbf E$'s units in newtons and in volts/metre; show they
agree. (AP-C)
\end{exercise}
\begin{exercise}[easy]
From a field-line map of two charges, one region has 3$\times$ the line
density of another: compare magnitudes. (F=ma)
\end{exercise}
\begin{exercise}[easy]
$+Q$ and $-Q$ sit $d$ apart. Where on the axis is $E=0$ --- if anywhere?
(AP-C)
\end{exercise}
\begin{exercise}[medium]
Derive the on-axis field of a uniformly charged rod of length $L$ at
distance $d$ from its end. (AP-C FR-style)
\end{exercise}
\begin{exercise}[medium]
A semicircle of charge $Q$, radius $R$: find $\mathbf E$ at the centre.
(AP-C)
\end{exercise}
\begin{exercise}[medium]
Three charges in a line: $+Q$, $-2Q$, $+Q$ equally spaced --- is the
middle one's equilibrium stable along the axis? (F=ma)
\end{exercise}
\begin{exercise}[medium]
A dipole of moment $p$ sits in a uniform field at angle $\theta$: find
the torque and the equilibrium direction. (AP-C)
\end{exercise}
\begin{exercise}[medium]
Design: plates $5$\,cm long must deflect $30$\,keV electrons by
$2$\,mm --- find $E$. (engineering)
\end{exercise}
\begin{exercise}[hard]
An arc subtending angle $\alpha$ has $\lambda=\lambda_0\sin\theta$:
set up the integral for $\mathbf E$ at the centre. (USAPhO-style)
\end{exercise}
\begin{exercise}[hard]
A dipole sits in a non-uniform field: argue it feels a net force toward
the stronger region and estimate it via $\mathbf p\cdot\nabla\mathbf E$.
(APhO-style)
\end{exercise}
\begin{exercise}[medium]
At what angle from the dipole axis does the far field's magnitude equal
the equatorial value at the same $r$? (F=ma)
\end{exercise}
\begin{exercise}[medium]
A point charge sits between two grounded conductors: explain, via field
lines, why its field is distorted toward the plates. (concept)
\end{exercise}
\begin{exercise}[medium]
Two equal charges $+Q$ at $\pm d/2$: find $\mathbf E$ on the axis
$x>0$ and its zero. (AP-C)
\end{exercise}
